%% toc: P15 170 170.5,170.6
% Proofs for repeated fixed-sample tests on accumulating data.
\section{Repeated fixed-sample tests on accumulating data: proofs}
\label{sec:annex-p15-w02b-p15-061-cumulative-looks}

Use the notation of \cref{def:p15-w02b-p15-061-procedure}. The variables
\(Y_i\) are independent and identically distributed with mean \(0\) and
variance \(1\). Since \(Y_i\) is an affine bijection of \(X_i\),
\(\sigma(Y_1,\ldots,Y_n)=\mathcal F_n\). For the transition kernels write
\[
 q_k(z,u)=\frac1{\sqrt{1-\rho_k^2}}\,
 \varphi\Bigl(\frac{u-\rho_kz}{\sqrt{1-\rho_k^2}}\Bigr),
\]
the \(N(\rho_kz,1-\rho_k^2)\) density in \(u\).

\begin{annexitem}{annex:p15-w02b-p15-061-look-dependence}{thm:p15-w02b-p15-061-look-dependence}{Covariances, the one-step representation and the joint density}
\begin{proof}
\emph{Covariances.} For \(j\leq k\),
\(\operatorname{Cov}_0(S_{n_j},S_{n_k})=\sum_{i\leq n_j}\sum_{l\leq n_k}
\operatorname{Cov}_0(Y_i,Y_l)\). Independence makes every term with
\(i\neq l\) zero, and the terms with \(i=l\leq n_j\) equal one, so the sum
is \(n_j\). Dividing by \(\sqrt{n_jn_k}\) gives
\(\operatorname{Cov}_0(Z_{n_j},Z_{n_k})=\sqrt{n_j/n_k}\).

\emph{One-step representation.} As \(0<n_k<n_{k+1}\), we have
\(\rho_k\in(0,1)\) and \(\sqrt{1-\rho_k^2}=
\sqrt{(n_{k+1}-n_k)/n_{k+1}}\). Hence
\[
 \rho_kZ_{n_k}+\sqrt{1-\rho_k^2}\,W_{k+1}
 =\frac{S_{n_k}}{\sqrt{n_{k+1}}}
 +\frac{S_{n_{k+1}}-S_{n_k}}{\sqrt{n_{k+1}}}
 =Z_{n_{k+1}} .
\]

\emph{Independence and moments.} \(Z_{n_1}\) is a function of
\(Y_1,\ldots,Y_{n_1}\) and \(W_{k+1}\) is a function of
\(Y_{n_k+1},\ldots,Y_{n_{k+1}}\). These blocks of indices are pairwise
disjoint, so \cref{thm:p15-independent-blocks} makes
\(Z_{n_1},W_2,W_3,\ldots\) independent, and makes \(W_{k+1}\) independent
of \(\sigma(Y_1,\ldots,Y_{n_k})=\mathcal F_{n_k}\). Each is a sum of
centered unit-variance independent terms divided by the square root of
their number, so it has mean \(0\) and variance \(1\).

\emph{Gaussian case.} If \(X_i\sim N(\mu_0,\sigma^2)\) then
\(Y_i\sim N(0,1)\). By Theorem~2.1.20 of
\cite[p.~52]{DurrettPTEAuthor2019} and induction on the number of
terms, a sum of \(m\) independent \(N(0,1)\) variables is \(N(0,m)\);
dividing by \(\sqrt m\) gives \(N(0,1)\). Thus \(Z_{n_1},W_2,\ldots,W_K\)
are independent standard normal, with joint density
\(\varphi(z_1)\prod_{k=2}^{K}\varphi(w_k)\) on \(\mathbb R^K\).
The map \((z_1,w_2,\ldots,w_K)\mapsto(z_1,\ldots,z_K)\) defined by
\(z_{k+1}=\rho_kz_k+\sqrt{1-\rho_k^2}\,w_{k+1}\) is a linear bijection of
\(\mathbb R^K\), with inverse
\(w_{k+1}=(z_{k+1}-\rho_kz_k)/\sqrt{1-\rho_k^2}\). The inverse has a
lower-triangular matrix with diagonal entries
\(1,(1-\rho_1^2)^{-1/2},\ldots,(1-\rho_{K-1}^2)^{-1/2}\), so its
Jacobian determinant is \(\prod_{k=1}^{K-1}(1-\rho_k^2)^{-1/2}\).
The change-of-variables formula for a linear bijection gives the density
\(f_K\) of \((Z_{n_1},\ldots,Z_{n_K})\) stated in the lemma. It is a
product of values of \(\varphi>0\) and positive constants, so it is
strictly positive everywhere. As a linear image of a standard normal
vector the vector is centered normal, with the covariances computed
above.
\end{proof}
\end{annexitem}

\begin{annexitem}{annex:p15-w02b-p15-061-finite-error}{thm:p15-w02b-p15-061-finite-error}{Finite-look error}
\begin{proof}
By \cref{thm:p15-w02b-p15-061-look-dependence}, applied to any initial
segment of the schedule, the first \(j\) looks have density \(f_j\), and
each \(Z_{n_k}\) is standard normal.

\emph{Part 1.} The rejection event has complement
\[
 \Bigl\{\max_{k\leq K}Z_{n_k}\leq c\Bigr\}
 =\bigl\{(Z_{n_1},\ldots,Z_{n_K})\in(-\infty,c]^K\bigr\},
\]
whose probability is \(\int_{(-\infty,c]^K}f_K\).

\emph{Part 2.} We have
\(f_K(z)=\varphi(z_1)\prod_{k=1}^{K-1}q_k(z_k,z_{k+1})\). We show by
backward induction on \(k\) that, for \(z_k\leq c\),
\[
 \int_{(-\infty,c]^{K-k}}\prod_{j=k}^{K-1}q_j(z_j,z_{j+1})\,
 dz_{k+1}\cdots dz_K=h_k(z_k).
\]
For \(k=K\) the product and the integral are empty and both sides equal
\(1=h_K\). If the identity holds for \(k+1\), integrate first over
\(z_{k+2},\ldots,z_K\) and then over \(z_{k+1}\); Tonelli's theorem
permits this order because the integrand is nonnegative. The inner
integral is \(h_{k+1}(z_{k+1})\) and the outer one is
\(\int_{-\infty}^cq_k(z_k,u)h_{k+1}(u)\,du=h_k(z_k)\). Multiplying the
case \(k=1\) by \(\varphi(z_1)\) and integrating over \(z_1\leq c\)
gives \(1-\alpha_K(c)=\int_{-\infty}^c\varphi(z)h_1(z)\,dz\) by part 1.

\emph{Part 3.} For \(K=2\),
\(h_1(z)=\int_{-\infty}^cq_1(z,u)\,du
=\Phi\bigl((c-\rho z)/\sqrt{1-\rho^2}\bigr)\), the distribution function
of \(N(\rho z,1-\rho^2)\) at \(c\). Writing \(\Phi=1-\bar\Phi\),
\[
 1-\alpha_2(c)=\int_{-\infty}^c\varphi(z)\,dz
 -\int_{-\infty}^c\varphi(z)\,\bar\Phi\Bigl(\frac{c-\rho z}{\sqrt{1-\rho^2}}\Bigr)dz ,
\]
and the first integral is \(1-\bar\Phi(c)\). Rearranging gives part 3.

\emph{Part 4.} Since \(Z_{n_1}\) is standard normal,
\(\alpha_1(c)=\bar\Phi(c)\). Let \(K\geq2\). The event
\(\{\max_{k\leq K}Z_{n_k}>c\}\) is the disjoint union of
\(\{\max_{k\leq K-1}Z_{n_k}>c\}\) and
\[
 D_K=\{Z_{n_1}\leq c,\ldots,Z_{n_{K-1}}\leq c,\ Z_{n_K}>c\}.
\]
Hence \(\alpha_K(c)-\alpha_{K-1}(c)=P_0(D_K)=
\int_{(-\infty,c]^{K-1}\times(c,\infty)}f_K\). The region contains a
box of positive volume and \(f_K>0\) everywhere, so this integral is
positive. Next, \(1-\alpha_K(c)=\int_{(-\infty,c]^K}f_K>0\) for the same
reason, so \(\alpha_K(c)<1\). Finally let \(A_k=\{Z_{n_k}>c\}\). By
subadditivity and inclusion--exclusion for two events,
\[
 P_0\Bigl(\bigcup_{k=1}^KA_k\Bigr)\leq P_0(A_1\cup A_2)
 +\sum_{k=3}^KP_0(A_k),
\]
and \(P_0(A_1\cup A_2)=P_0(A_1)+P_0(A_2)-P_0(A_1\cap A_2)\). As each
\(P_0(A_k)=\bar\Phi(c)\), the right side is
\(K\bar\Phi(c)-P_0(A_1\cap A_2)\),
and \(P_0(A_1\cap A_2)=\int_{(c,\infty)^2}f_2>0\). So
\(\alpha_K(c)<K\bar\Phi(c)\). With \(c=\bar\Phi^{-1}(\alpha)\) each
preselected look has level \(P_0(Z_{n_k}>c)=\bar\Phi(c)=\alpha\), while
\(\alpha_K(c)\geq\alpha_2(c)>\alpha_1(c)=\alpha\) for \(K\geq2\).

\emph{Part 5.} Let \(R\) be the rejection event of a rule that inspects
a data-dependent subset of the looks \(n_1,\ldots,n_K\) and rejects at an
inspected look with \(Z_{n_k}>c\). On \(R\) some \(k\leq K\) has
\(Z_{n_k}>c\), so \(R\subseteq\{\max_{k\leq K}Z_{n_k}>c\}\) and
\(P_0(R)\leq\alpha_K(c)\) by monotonicity.
\end{proof}
\end{annexitem}

\begin{annexitem}{annex:p15-w02b-p15-061-association}{thm:p15-w02b-p15-061-association}{Monotone functions of independent variables}
\begin{proof}
Call a function \(h\colon\mathbb R^m\to[0,1]\) \emph{decreasing-type} if
it is Borel and nonincreasing in each coordinate when the other
coordinates are held fixed. Let \(\mu_i\) be the law of \(U_i\) and
\(\mu^{(m)}=\mu_1\times\cdots\times\mu_m\). By
\cref{thm:p15-independent-product-law} the law of \(U\) on
\(\mathcal B(\mathbb R^m)\) is \(\mu^{(m)}\), so
\cref{thm:p15-expectation-law-integral} gives
\(E[h(U)]=\int h\,d\mu^{(m)}\) for every bounded Borel \(h\). It
therefore suffices to prove, by induction on \(m\geq1\), the following
claim.
\begin{quote}
\((C_m)\) For all Borel probability measures \(\mu_1,\ldots,\mu_m\) on
\(\mathbb R\) and all decreasing-type \(f,g\) on \(\mathbb R^m\),
\(\int fg\,d\mu^{(m)}\geq\int f\,d\mu^{(m)}\int g\,d\mu^{(m)}\).
\end{quote}
Every integrand below is Borel with absolute value at most \(1\) and is
integrated against a probability measure, so every integral is finite and
linearity (\cref{thm:p15-signed-integral-laws}) applies.

\emph{The case \(m=1\).} For bounded Borel \(f,g\colon\mathbb R\to\mathbb R\)
and a Borel probability measure \(\mu\) on \(\mathbb R\), put
\(D(u,v)=(f(u)-f(v))(g(u)-g(v))\), a bounded Borel function on
\(\mathbb R^2\). Expanding,
\(D(u,v)=f(u)g(u)-f(u)g(v)-f(v)g(u)+f(v)g(v)\). Each of the four terms
is bounded and Borel, so \cref{thm:p15-fubini} evaluates its integral
against the probability measure \(\mu\times\mu\) as an iterated integral:
the first and last give \(\int fg\,d\mu\), and the middle two give
\(\int f\,d\mu\int g\,d\mu\). Hence
\[
 \int_{\mathbb R^2}D\,d(\mu\times\mu)
 =2\Bigl(\int fg\,d\mu-\int f\,d\mu\int g\,d\mu\Bigr).
 \tag{\(\ast\)}
\]
This identity uses no monotonicity. Now let \(f,g\) be decreasing-type.
If \(u\leq v\), then \(f(u)-f(v)\geq0\) and \(g(u)-g(v)\geq0\); if
\(u\geq v\), both differences are at most \(0\). In either case
\(D(u,v)\geq0\), so the left side of
\((\ast)\) is nonnegative, which is
\((C_1)\).

\emph{The step from \(m-1\) to \(m\).} Let \(m\geq2\) and assume
\((C_{m-1})\). Write points of \(\mathbb R^m\) as \((u',t)\) with
\(u'\in\mathbb R^{m-1}\) and \(t\in\mathbb R\), and put
\(\nu=\mu_1\times\cdots\times\mu_{m-1}\). By
\cref{thm:p15-product-sections},
\(\mathcal B(\mathbb R^m)=\mathcal B(\mathbb R^{m-1})\times\mathcal B(\mathbb R)\),
and by \cref{thm:p15-finite-product-measures}
\(\mu^{(m)}=\nu\times\mu_m\) under this identification. Define
\begin{gather*}
 F(u')=\int f(u',t)\,\mu_m(dt),\qquad
 G(u')=\int g(u',t)\,\mu_m(dt),\\
 H(u')=\int f(u',t)g(u',t)\,\mu_m(dt).
\end{gather*}
The sections \(t\mapsto f(u',t)\) and \(t\mapsto g(u',t)\) are Borel by
\cref{thm:p15-product-sections}, and \cref{thm:p15-tonelli}, applied to
the nonnegative functions \(f\), \(g\) and \(fg\), shows that \(F\),
\(G\) and \(H\) are Borel functions of \(u'\) with
\[
 \int f\,d\mu^{(m)}=\int F\,d\nu,\qquad
 \int g\,d\mu^{(m)}=\int G\,d\nu,\qquad
 \int fg\,d\mu^{(m)}=\int H\,d\nu .
\]

(a) For fixed \(u'\), the sections \(t\mapsto f(u',t)\) and
\(t\mapsto g(u',t)\) take values in \([0,1]\) and are nonincreasing in
\(t\), since \(t\) is one coordinate of \(f\) and \(g\). So they are
decreasing-type on \(\mathbb R\), and \((C_1)\) with the measure
\(\mu_m\) gives \(H(u')\geq F(u')G(u')\) for every \(u'\).

(b) \(F\) and \(G\) take values in \([0,1]\), as integrals of
\([0,1]\)-valued functions against a probability measure. Let \(u'\) and
\(w'\) differ only in coordinate \(i\), with \(u'_i\leq w'_i\). For every
\(t\) the points \((u',t)\) and \((w',t)\) differ only in coordinate
\(i\), so \(f(u',t)\geq f(w',t)\) and \(g(u',t)\geq g(w',t)\);
integrating against \(\mu_m\) gives \(F(u')\geq F(w')\) and
\(G(u')\geq G(w')\). Thus \(F\) and \(G\) are decreasing-type on
\(\mathbb R^{m-1}\), and \((C_{m-1})\) with the measures
\(\mu_1,\ldots,\mu_{m-1}\) gives
\(\int FG\,d\nu\geq\int F\,d\nu\int G\,d\nu\).

Integrating the pointwise inequality of (a) against \(\nu\) and then
using (b),
\[
 \int fg\,d\mu^{(m)}=\int H\,d\nu\geq\int FG\,d\nu
 \geq\int F\,d\nu\int G\,d\nu
 =\int f\,d\mu^{(m)}\int g\,d\mu^{(m)} ,
\]
which is \((C_m)\). By induction \((C_m)\) holds for every \(m\), and the
lemma follows by the reduction in the first paragraph.
\end{proof}
\end{annexitem}

\begin{annexitem}{annex:p15-w02b-p15-061-versus-independent}{thm:p15-w02b-p15-061-versus-independent}{Overlapping looks compared with independent looks}
\begin{proof}
\emph{Independent block sums.} Put \(n_0=0\) and, for \(1\leq k\leq K\),
\[
 V_k=S_{n_k}-S_{n_{k-1}}=\sum_{i=n_{k-1}+1}^{n_k}Y_i ,
\]
so that \(S_{n_k}=V_1+\cdots+V_k\). Each \(V_k\) is a Borel function (a
finite sum) of the block \(Y_{n_{k-1}+1},\ldots,Y_{n_k}\). These blocks
are pairwise disjoint and the \(Y_i\) are independent, so
\cref{thm:p15-independent-blocks} makes \(V_1,\ldots,V_K\) independent
real random variables. Put \(V=(V_1,\ldots,V_K)\).

\emph{Acceptance indicators.} For \(1\leq k\leq K\) let
\(a_k\colon\mathbb R^K\to[0,1]\) be the indicator of the closed set
\(\{v:v_1+\cdots+v_k\leq c\sqrt{n_k}\}\); it is Borel. Increasing a
coordinate \(v_j\) with \(j\leq k\) increases \(v_1+\cdots+v_k\) and can
only change the value of \(a_k\) from \(1\) to \(0\), and coordinates
\(v_j\) with \(j>k\) do not enter, so \(a_k\) is nonincreasing in each
coordinate. Since \(Z_{n_k}\leq c\) exactly when
\(S_{n_k}\leq c\sqrt{n_k}\), we have \(a_k(V)=\mathbf 1_{G_k}\) with
\(G_k=\{Z_{n_k}\leq c\}\).
If \(p,q\colon\mathbb R^K\to[0,1]\) are Borel and nonincreasing in each
coordinate, so is \(pq\): when \(u\) and \(w\) differ only in one
coordinate and \(u\) is smaller there,
\(p(u)q(u)\geq p(w)q(u)\geq p(w)q(w)\), because \(q(u)\geq0\) and
\(p(w)\geq0\). Hence \(b_j=a_1\cdots a_j\) is Borel, \([0,1]\)-valued and
nonincreasing in each coordinate, and
\(b_j(V)=\mathbf 1_{G_1\cap\cdots\cap G_j}\).

\emph{The general inequality.} For \(2\leq j\leq K\),
\cref{thm:p15-w02b-p15-061-association} with \(m=K\), \(U=V\),
\(f=b_{j-1}\) and \(g=a_j\) gives
\[
 P_0(G_1\cap\cdots\cap G_j)=E\bigl[b_{j-1}(V)a_j(V)\bigr]
 \geq P_0(G_1\cap\cdots\cap G_{j-1})\,P_0(G_j).
\]
Applying this for \(j=K,K-1,\ldots,3\) and multiplying the inequalities
by the nonnegative factors \(P_0(G_k)\) as they accumulate gives, for
\(K\geq2\),
\[
 P_0\Bigl(\bigcap_{k=1}^KG_k\Bigr)\geq
 P_0(G_1\cap G_2)\prod_{k=3}^KP_0(G_k)
 \tag{\(\ast\ast\)}
\]
where the product is empty when \(K=2\). The case \(j=2\) then gives
\[
 P_0\Bigl(\bigcap_{k=1}^KG_k\Bigr)\geq\prod_{k=1}^KP_0(G_k).
\]
For \(K=1\) the two sides are equal. Since
\(\bigcap_{k\leq K}G_k=\{\max_{k\leq K}Z_{n_k}\leq c\}\), the left side
is \(1-\alpha_K(c)\), which proves the displayed inequality.

\emph{Strict inequality in the Gaussian null model.} Let \(K\geq2\). By
\cref{thm:p15-w02b-p15-061-look-dependence} each \(Z_{n_k}\) is standard
normal, so \(P_0(G_k)=\Phi(c)>0\). By
\((\ast\ast)\) it suffices to show
\(P_0(G_1\cap G_2)>\Phi(c)^2\): multiplying by
\(\Phi(c)^{K-2}>0\) then gives \(1-\alpha_K(c)>\Phi(c)^K\).

Write \(Z=Z_{n_1}\), \(W=W_2\), \(\rho=\rho_1\in(0,1)\) and
\(s=\sqrt{1-\rho^2}>0\). By
\cref{thm:p15-w02b-p15-061-look-dependence}, \(Z\) and \(W\) are
independent standard normal and \(Z_{n_2}=\rho Z+sW\). Let \(\gamma\) be
the standard normal law, \(a(z)=\mathbf 1_{\{z\leq c\}}\) and
\[
 h(z)=P_0(\rho z+sW\leq c)=\Phi\Bigl(\frac{c-\rho z}{s}\Bigr)
 \qquad(z\in\mathbb R).
\]
As \(\varphi>0\), \(\Phi\) is strictly increasing; since \(\rho>0\) and
\(s>0\), \(h\) is strictly decreasing. Applying
\cref{thm:p15-independent-expectations} to the bounded Borel functions
\((z,w)\mapsto a(z)\mathbf 1_{\{\rho z+sw\leq c\}}\) and
\((z,w)\mapsto\mathbf 1_{\{\rho z+sw\leq c\}}\) of \((Z,W)\), and
integrating over \(w\) first,
\[
 P_0(G_1\cap G_2)=\int ah\,d\gamma,\qquad
 P_0(G_2)=\int h\,d\gamma,\qquad
 P_0(G_1)=\int a\,d\gamma .
\]
Both \(P_0(G_1)\) and \(P_0(G_2)\) equal \(\Phi(c)\). The identity
\((\ast)\) in the proof of
\cref{thm:p15-w02b-p15-061-association}, with \(\mu=\gamma\), \(f=a\)
and \(g=h\), gives
\[
 P_0(G_1\cap G_2)-\Phi(c)^2
 =\frac12\int_{\mathbb R^2}\bigl(a(u)-a(v)\bigr)\bigl(h(u)-h(v)\bigr)
 \varphi(u)\varphi(v)\,du\,dv ,
\]
where \(\gamma\times\gamma\) has been written with its density
\(\varphi(u)\varphi(v)\) with respect to Lebesgue measure on
\(\mathbb R^2\), by \cref{thm:p15-tonelli} applied to the nonnegative
parts of the integrand. Both \(a\) and \(h\) are nonincreasing, so the
integrand is nonnegative, as in the case \(m=1\) above. On the rectangle
\(R=(c-1,c]\times(c,c+1]\), which has Lebesgue measure \(1\), we have
\(u<v\), \(a(u)-a(v)=1\), \(h(u)-h(v)>0\) and
\(\varphi(u)\varphi(v)>0\), so the integrand is strictly positive on
\(R\). A nonnegative function whose integral is zero vanishes almost
everywhere (\cref{thm:p15-nonnegative-integral-laws}), so the integral is
strictly positive and \(P_0(G_1\cap G_2)>\Phi(c)^2\).

\emph{The bounds on \(\alpha_K(c)\).} From the strict inequality,
\(\alpha_K(c)<1-\Phi(c)^K=1-(1-\bar\Phi(c))^K\). Part 4 of
\cref{thm:p15-w02b-p15-061-finite-error} gives
\(\alpha_1(c)<\alpha_2(c)<\cdots<\alpha_K(c)\) and
\(\alpha_1(c)=\bar\Phi(c)\), so \(\bar\Phi(c)<\alpha_K(c)\). The function
\(\bar\Phi\) is a continuous strictly decreasing bijection of
\(\mathbb R\) onto \((0,1)\), so for \(0<\alpha<1\) the threshold
\(c=\bar\Phi^{-1}(\alpha)\) exists and \(\bar\Phi(c)=\alpha\); substituting
gives \(\alpha<\alpha_K(c)<1-(1-\alpha)^K\). The right side is the value
\(1-(1-\alpha)^M\) of \cref{thm:p15-se-repeated-fixed} with \(M=K\).
\end{proof}
\end{annexitem}

\begin{annexitem}{annex:p15-w02b-p15-061-unbounded}{thm:p15-w02b-p15-061-unbounded}{The upper limit of the looks is a tail variable equal to infinity}
\begin{proof}
Let \(L=\limsup_{k\to\infty}Z_{n_k}\), an extended-real random variable.
Since the \(n_k\) are strictly increasing integers, \(n_k\geq k\) and
\(n_k\to\infty\).

\emph{Step 1: \(L\) is a tail variable.} Fix \(m\). For \(n_k>m\),
\[
 Z_{n_k}=\frac{S_m}{\sqrt{n_k}}+\frac{S_{n_k}-S_m}{\sqrt{n_k}} .
\]
At every sample point \(S_m\) is a real number, so the first term tends
to \(0\) as \(k\to\infty\). Therefore
\(L=\limsup_k(S_{n_k}-S_m)/\sqrt{n_k}\) everywhere. Each
\((S_{n_k}-S_m)/\sqrt{n_k}\) is a function of
\(Y_{m+1},Y_{m+2},\ldots\), hence measurable with respect to
\(\sigma(X_{m+1},X_{m+2},\ldots)\), and an
upper limit of measurable extended-real functions is measurable. As
\(m\) was arbitrary, \(L\) is measurable with respect to the tail
sigma-algebra \(\mathcal T=\bigcap_m\sigma(X_{m+1},X_{m+2},\ldots)\);
compare Example~2.5.2 of \cite[p.~81]{DurrettPTEAuthor2019}. For every
real \(x\), \(\{L\geq x\}\in\mathcal T\), and
\cref{thm:p15-kolmogorov-tail-zero-one} (Theorem~2.5.3 of
\cite[p.~81]{DurrettPTEAuthor2019}) gives
\(P_0(L\geq x)\in\{0,1\}\).

\emph{Step 2: \(P_0(L\geq x)>0\).} Put \(A_k=\{Z_{n_k}>x\}\). The
variables \(Y_i\) are independent and identically distributed with mean
\(0\) and variance \(1\), so by \cref{thm:p15-clt-iid} (Theorem~3.4.1 of
\cite[p.~144]{DurrettPTEAuthor2019}) \(Z_n\) converges in distribution to
a standard normal variable. Since \(\Phi\) is continuous at \(x\),
\(P_0(Z_n>x)\to\bar\Phi(x)\), and along the subsequence
\(P_0(A_k)\to\bar\Phi(x)>0\). The events
\(B_m=\bigcup_{k\geq m}A_k\) decrease in \(m\) and
\(P_0(B_m)\geq\sup_{k\geq m}P_0(A_k)\). By continuity from above,
\[
 P_0(A_k\text{ for infinitely many }k)=\lim_mP_0(B_m)
 \geq\lim_m\sup_{k\geq m}P_0(A_k)=\bar\Phi(x)>0 .
\]
If \(A_k\) occurs for infinitely many \(k\), then
\(\sup_{k\geq m}Z_{n_k}>x\) for every \(m\), so \(L\geq x\). Hence
\(P_0(L\geq x)>0\), and by step 1 \(P_0(L\geq x)=1\).

\emph{Step 3: conclusion.} The event \(\{L=+\infty\}\) is the
intersection of the probability-one events \(\{L\geq j\}\),
\(j=1,2,\ldots\), so it has probability one. On it, for every \(m\),
\(\sup_{k\geq m}Z_{n_k}=+\infty>c\), so some \(k\geq m\) has
\(Z_{n_k}>c\); thus \(Z_{n_k}>c\) for infinitely many \(k\). That event
is contained in \(\{T<\infty\}\), so \(\alpha_\infty(c)=1\). The events
\(\{T\leq K\}\) increase with \(K\) to \(\{T<\infty\}\), so continuity
from below gives \(\alpha_K(c)\uparrow1\).
\end{proof}
\end{annexitem}

\begin{annexitem}{annex:p15-w02b-p15-061-two-sided}{thm:p15-w02b-p15-061-two-sided}{The two-sided rule}
\begin{proof}
In the Gaussian null model \(Z_{n_1}\) is standard normal, so by the
symmetry of \(\varphi\)
\[
 \alpha^{\pm}_1(c)=P_0(|Z_{n_1}|>c)=2\bar\Phi(c).
\] The complement of \(\{\max_{k\leq K}|Z_{n_k}|>c\}\) is
\(\{(Z_{n_1},\ldots,Z_{n_K})\in[-c,c]^K\}\), which gives the integral.
The recursion follows from the induction in part 2 of
\cref{thm:p15-w02b-p15-061-finite-error}, carried out with every
integral over \([-c,c]\) instead of \((-\infty,c]\); that induction used
only the product form of \(f_K\) and Tonelli's theorem, not the shape of
the interval. For \(K\geq2\) the increment
\(\alpha^{\pm}_K(c)-\alpha^{\pm}_{K-1}(c)\) is the integral of \(f_K\)
over \([-c,c]^{K-1}\times(\mathbb R\setminus[-c,c])\), and
\(1-\alpha^{\pm}_K(c)\) is its integral over \([-c,c]^K\). Because
\(c>0\), both regions contain boxes of positive volume, and \(f_K>0\), so
both are positive. With \(E_k=\{|Z_{n_k}|>c\}\), the same subadditivity
estimate as in part 4 of that proof gives
\(\alpha^{\pm}_K(c)\leq2K\bar\Phi(c)-P_0(E_1\cap E_2)\), and
\(P_0(E_1\cap E_2)\) is the integral of \(f_2\) over
\((\mathbb R\setminus[-c,c])^2\), which is positive. For the unbounded
statement, \(\{Z_{n_k}>c\}\subseteq\{|Z_{n_k}|>c\}\) for every \(k\), so
the two-sided rule rejects whenever the one-sided rule does, and
\(\alpha^{\pm}_\infty(c)\geq\alpha_\infty(c)=1\) by
\cref{thm:p15-w02b-p15-061-unbounded}.
\end{proof}
\end{annexitem}

\begin{annexitem}{annex:p15-w02b-p15-061-calibrated}{thm:p15-w02b-p15-061-calibrated}{Continuity and strict monotonicity of the finite-look error}
\begin{proof}
Let \(M=\max_{k\leq K}Z_{n_k}\) and \(G(c)=P_0(M\leq c)=1-\alpha_K(c)\).
As the distribution function of a real random variable, \(G\) is
nondecreasing and right-continuous, with limits \(0\) at \(-\infty\) and
\(1\) at \(+\infty\). Its jump at \(c\) is
\(P_0(M=c)\leq\sum_{k\leq K}P_0(Z_{n_k}=c)=0\), because each
\(Z_{n_k}\) is standard normal; so \(G\) is continuous. For \(c<c'\),
if every \(Z_{n_k}\) lies in \((c,c']\) then \(c<M\leq c'\), so
\[
 G(c')-G(c)=P_0(c<M\leq c')\geq\int_{(c,c']^K}f_K>0
\]
by \cref{thm:p15-w02b-p15-061-look-dependence}. Thus \(\alpha_K=1-G\) is
continuous and strictly decreasing from \(1\) to \(0\). For
\(\alpha\in(0,1)\) the intermediate value theorem gives some \(c_K\) with
\(\alpha_K(c_K)=\alpha\), and strict monotonicity makes it unique. Since
\(\alpha_1(c)=\bar\Phi(c)\), \(c_1=\bar\Phi^{-1}(\alpha)\). For
\(K\geq2\), part 4 of \cref{thm:p15-w02b-p15-061-finite-error} at the
point \(c_{K-1}\) gives
\(\alpha_K(c_{K-1})>\alpha_{K-1}(c_{K-1})=\alpha=\alpha_K(c_K)\); as
\(\alpha_K\) is strictly decreasing, \(c_{K-1}<c_K\). The first-crossing
rule with threshold \(c_K\) has error \(\alpha_K(c_K)=\alpha\) by
definition. The Gaussian null model satisfies the hypotheses of
\cref{thm:p15-w02b-p15-061-unbounded}, so with an infinite schedule
\(\alpha_\infty(c)=1\) for every real \(c\).
\end{proof}
\end{annexitem}
